\section{制御回転による振幅の変換}
\label{sec:controlled-rotation}

値レジスタの計算基底を$\{\ket z\}$とし，
実数値関数$g$がすべての$z$について$|g(z)|\leq1$を満たすとする．
値$z$に応じて補助量子ビットを回転する演算$W_g$を
\begin{equation}
 W_g\ket z\ket0
 =\ket z\left(
   \sqrt{1-g(z)^2}\ket0+g(z)\ket1
 \right)
 \label{eq:controlled-rotation-general}
\end{equation}
と定める．

回転角を$\theta_z=2\arcsin g(z)$とおけば，
式\eqref{eq:ry-on-zero}より補助量子ビットへの$R_y(\theta_z)$によって
式\eqref{eq:controlled-rotation-general}が得られる．
値レジスタ全体では
\begin{equation}
 W_g=\sum_z\ket z\bra z\otimes R_y(2\arcsin g(z))
 \label{eq:controlled-rotation-block}
\end{equation}
と書ける．

\begin{prop}
演算$W_g$はユニタリである．また，規格化された入力
$\sum_z\alpha_z\ket z\ket0$に$W_g$を作用させて補助量子ビットを測定すると，
結果$1$を得る確率は
\begin{equation}
 p_1=\sum_z|\alpha_z|^2g(z)^2
 \label{eq:controlled-rotation-success-probability}
\end{equation}
である．成功確率が$p_1>0$のとき，値レジスタの条件付き状態は
\begin{equation}
 \frac1{\sqrt{p_1}}\sum_z\alpha_zg(z)\ket z
 \label{eq:controlled-rotation-conditional-state}
\end{equation}
である．
\end{prop}
\begin{proof}
式\eqref{eq:controlled-rotation-block}は，互いに直交する部分空間
$\operatorname{span}\{\ket z\}\otimes\mathbb C^2$ごとに
ユニタリ$R_y(2\arcsin g(z))$を作用させるブロック対角演算である．実際，
\begin{align*}
 W_g^\dagger W_g
 &=\sum_{z,z'}\ket z\braket{z|z'}\bra{z'}
   \otimes R_y(2\arcsin g(z))^\dagger R_y(2\arcsin g(z'))\\
 &=\sum_z\ket z\bra z\otimes I=I
\end{align*}
である．入力へ作用させた状態は
\[
 \sum_z\alpha_z\sqrt{1-g(z)^2}\ket z\ket0
 +\sum_z\alpha_zg(z)\ket z\ket1
\]
となる．第2項のノルムの二乗を計算基底の正規直交性で求めると
式\eqref{eq:controlled-rotation-success-probability}を得る．
射影後の第2項を$\sqrt{p_1}$で割れば
式\eqref{eq:controlled-rotation-conditional-state}となる．
\end{proof}

HHLでは，$z$が固有値の近似値$\widetilde\lambda$を表し，
$g(\widetilde\lambda)=C/\widetilde\lambda$と選ぶ．
定数$C$の条件，固有値の符号，有限精度による誤差は第4章で定める．
本節では，角度を計算する可逆回路の構成までは扱わず，
式\eqref{eq:controlled-rotation-block}の演算を必要な精度で実行できるものとする．
